\section{Relevanz des Projekts} \todo{Gegenlesen}
Vor diesem Hintergrund gewinnt die Entwicklung einer eigenen Wortuhr besondere Relevanz. 
Ziel ist es, bestehende Schwächen der aktuellen Marktsegmente gezielt zu adressieren und eine Lösung zu entwickeln, die gestalterische und funktionale Anforderungen integriert. 
Die entwickelte Wortuhr soll folgenden, grundlegenden und zusätzliche Funktionen aufweisen und darstellen:

\begin{enumerate}
	\item Kontinuierliche und unterbrechungsfreie wortbasierte Zeitanzeige von Stunde, Minute und Sekunde
	\item Wortbasierte Anzeige der Außentemperatur auf Basis definierter Schwellenwerte
	\item Wortbasierte Darstellung von Wetterinformationen
	\item Symbolische Darstellung der Tageszeit
	\item Symbolische Darstellung der Mondphase
	\item Symbolische Darstellung der Raumfeuchtigkeit
	\item Symbolische Darstellung der Innenraumtemperatur
\end{enumerate}

Die Bewertung in Tabelle~\ref{tab:bewertung_wordclock} zeigt, dass bestehende Systeme entweder in einzelnen Kategorien hohe Ausprägungen erreichen oder deutliche Defizite aufweisen, jedoch keine ausgewogene Gesamtlösung darstellen. 
Premiumprodukte erzielen hohe Werte im Bereich Design und Benutzerfreundlichkeit, weisen jedoch Defizite im Funktionsumfang und in der Erweiterbarkeit auf. 
DIY-Systeme hingegen bieten eine hohe Funktionalität und Erweiterbarkeit, erfüllen jedoch häufig nicht die Anforderungen an Design und Benutzerfreundlichkeit. 

\begin{table}[h]
\centering
\begin{tabular}{|c|c|c|c|c|c|c|}
\hline
\textbf{Produkt} & \textbf{Design} & \textbf{Funktion} & \textbf{Usability} & \textbf{Erweiterbarkeit} & \textbf{Preis-Leistung} & \textbf{Gesamt} \\
\hline
QLOCKTWO & 5 & 2 & 5 & 1 & 2 & 3,05 \\
\hline
Jupiter & 3 & 3 & 4 & 2 & 3 & 3,05 \\
\hline
Massenmarkt & 2 & 1 & 3 & 1 & 4 & 2,10 \\
\hline
Eigene Wortuhr & 4 & 5 & 4 & 5 & 4 & 4,45 \\
\hline
\end{tabular}
\caption{Bewertung verschiedener WordClock-Konzepte}
\label{tab:bewertung_wordclock}
\end{table} \todo{Tabelle überarbeiten und page matching}

Vor diesem Hintergrund wird im Rahmen dieser Arbeit das Ziel verfolgt, eine Wortuhr zu entwickeln, die in allen betrachteten Kategorien hohe Bewertungen erreicht. 
Hierdurch soll eine ausgewogene und leistungsfähige Gesamtlösung realisiert werden.
Damit positioniert sich die entwickelte Lösung zwischen den bestehenden Marktsegmenten und adressiert eine bislang nur unzureichend besetzte Nische. 
Sie stellt nicht nur eine technische Erweiterung bestehender Systeme dar, sondern verfolgt den Ansatz, die Wortuhr von einem Designobjekt zu einem Informationssystem weiterzuentwickeln. 
Die Kombination aus gestalterischer Qualität und funktionaler Erweiterbarkeit begründet die Relevanz der eigenen Entwicklung sowohl aus technischer als auch aus marktbezogener Perspektive.
